\documentclass[../main.tex]{subfiles}

\begin{document}

\chapter{Background}\label{ch:background}

This chapter explains key concepts and terminalogy relevant to this thesis.
It establishes the theoretical and technical foundation necessary for understanding the development of 
the proposed monitoring tool.
The contents are structured as follows: Section~\ref{sec:fim} introduces the concept of Federated Identity 
Management and describes the roles of the Identity Provider and Service Provider, including their interaction 
during the authentication process. 
Section~\ref{sec:shibboleth} introduces Shibboleth as an open-source implementation of SAML, with a 
particular focus on the Shibboleth IdP, which is highly relevant to this thesis. 
Finally, Section~\ref{sec:saml} provides an overview of the SAML framework and its role in 
federated authentication.

\section{Technical Foundation of Federated Identity Management}\label{sec:fim}

As already mentioned in the introduction, FIM allows users to access all their resources they 
need by using a single digital identity. 
It provides \acrfull{sso} functionalities across all applications of an organization. 
Through the properties granted by incorporating SSO into the federation, it eliminated the need 
to create, remember and manage multiple usernames and passwords for users and the organization. 
By centralising identity management, the number of systems holding credentials is reduced. 
Fewer accounts and identity management systems means fewer vectors to attack and credential of a user to steal. 
It also eases the identity management by centralizing it on a server. 
This centralisation introduces new risks. 
FIM relies on a foundation of trust between organizations and this trust is heavily reliant on 
the security of the IdP. 
The centralization by incorporating an IdP means it creates a single point of trust, 
or in different words, it concentrates the risk onto a single server. 

A compromise of a primary IdP could potentially grant broad unauthorised access across the 
entire federation. 
This characteristic requires both the IdP and SP to maintain a high level of security. 
The policies defined for each SP in the IdP must accurately translate to the rules 
at the corresponding SP. 
The configuration of the IdP and that of each SP must remain consistent with one another. 
Certificates, keys and protocol settings are maintained separately on both sides, and a mismatch between 
them can weaken the security of a transaction without either party noticing.
To better understand the interaction within a federated identity system, the roles and responsibilities of 
the IdP and SP are described in the following.

A federation consists of two kinds of participants. 
The IdP authenticates users and issues statements about them. 
The SP offers the resource a user wants to reach and relies on those statements instead of authenticating 
the user itself. 
Neither role is meaningful without the other, and their interaction follows a fixed pattern. 

When a user requests a protected resource, the SP does not check credentials. 
It determines which IdP is responsible for the user and redirects the browser there together 
with an authentication request. 
The IdP authenticates the user by whatever means the deployment provides, and returns a signed statement 
asserting that the authentication took place and describing the user by a set of attributes. 
The SP verifies the signature, evaluates the statement and grants or denies access on that basis. 
The credentials themselves never leave the IdP.

This division has a consequence for how a federation is configured. 
Because the IdP serves many SPs, it does not treat them uniformly. 
It maintains a separate configuration for each partner, specifying for instance whether the statement 
is to be encrypted, which authentication methods are permitted, and which attributes are released. 
The SP maintains its own configuration in turn, stating what it expects to receive and what 
it will accept.

Both sides therefore hold configuration that describes the same relationship, and both are 
maintained independently. 
What binds them together is the metadata each party publishes about itself: the endpoints at which 
it can be reached, and the certificates by which its statements can be verified. 
Trust within a federation rests on these metadata being current and correctly exchanged. 
Whether the two configurations remain consistent with one another, however, is not something the 
protocol itself verifies.

\section{Shibboleth}\label{sec:shibboleth}

\section{SAML}\label{sec:saml}


\end{document}
